\begin{frame}
	\frametitle{Ответы на замечания Крайнова~В.П.}
	\textbf{Замечание}: В главе~3 при исследовании эволюции кривых фазового синхронизма при фотоиндуцированной ГВГ в волокне наблюдается заметное расхождение экспериментальных данных кинетики эффективности ГВГ с результатами моделирования на начальной стадии записи решётки (t < 39 ч). С чем может быть связано это расхождение?
	
	\textbf{Замечание}: При обработке температурных кинетик в главе~4 предложенная итерационная процедура коррекции требует знания скорости теплоотдачи. Какова сходимость этой итерационной процедуры и насколько устойчивы полученные результаты к выбору конкретного вида функции $T_{a}(t)$? 
	
	\textbf{Ответ}: Процедура сходится за 2–3 итерации: изменение $\alpha$ на последних итерациях не превышает десятых долей процента. Устойчивость к виду $T_{a}(t)$ обеспечивается Фурье-разложением: поправка вычисляется для каждой гармоники независимо, поэтому результат не зависит от конкретной формы дрейфа.
\end{frame}